%% =============================================================
%%  chapters/05_data_acquisition.tex
%% =============================================================
\chapter{Dataset Acquisition and Synthetic Data Generation}
\label{ch:data-acquisition}

\section{Ethics and Data Privacy Considerations}
\label{sec:ethics-privacy}



\section{ Synthetic Pressure Sensor Data Generation}
\label{sec:synthetic-pressure-data}
 %%TODO : whats the approach to generate synthetic data, justification for choosing this one.
 %% what are its features, limitation, and considerations are made for this task.


\subsection{Sensor Hardware and Seat Setup Specifications}
\label{sec:sensor-hardware-seat-setup}

%% TODO: Insert exact sensor grid dimensions, sampling rate, and
%% calibration procedure used in the laboratory setup. seat setup images 

\subsection{Subjects and Experimental Protocol}
\label{sec:subjects-protocol}

%% TODO: Report exact number of subjects, age/weight/height distribution,
%% and the full list of seating configurations recorded.


\subsection{Data Statistics and Limitations of setup}
\label{sec:data-statistics-limitations}
The resulting real pressure dataset provides ground-truth labels for
occupancy, passenger class, weight, and coarse pose for each recorded
configuration, but is limited in overall subject diversity and recording
duration due to the practical constraints of laboratory-based data
collection. These limitations directly motivate the synthetic data
generation pipeline described in Section~\ref{sec:synthetic-heatmap-generation}.

%% TODO : state the limitations of the simulation environment.
%% TODO : 
%%------------------------------------------------------------------------------------------------------------------------
\section{Camera Data Collection}
\label{sec:camera-data-collection}

\subsection{Camera Hardware and Mounting}
\label{sec:camera-hardware-mounting}
Camera recordings were captured using an infrared-sensitive camera module
mounted in the vehicle headliner, providing a field of view that covers
the target seating positions under both daylight and night-time
illumination, following mounting conventions similar to those reported in
\cite{schmidt2020driver}.
%% TODO: Specify exact camera module, lens field of view, and mounting
%% position/angle used in the data collection rig.

\subsection{Recording Protocol}
\label{sec:camera-recording-protocol}
Camera frames were recorded continuously and time-stamped alongside the
pressure sensor stream for every subject session described in
Section~\ref{sec:subjects-protocol}, ensuring that each seating
configuration has a corresponding synchronised camera-pressure sample
pair available for the synchronisation pipeline in
Chapter~\ref{ch:preprocessing-sync}.
%% TODO: State recorded frame rate and resolution, and whether IR
%% illumination was active during all sessions.
%%-------------------------------------------------------------------------------------------------------------

\section{Synthetic Pressure Heatmap Generation}
\label{sec:synthetic-heatmap-generation}

\subsection{Human Body Simulation Model}
\label{sec:human-body-simulation-model}
Synthetic pressure heatmaps are generated by simulating a parametric
human body model with sampled shape and pose parameters seated on a
physics-based model of the seat geometry, following the simulation
approach proposed by \cite{ivanov2021synthetic}. Contact forces computed
by the physics simulation are projected onto a virtual sensor grid
matching the real sensor's spatial resolution to produce a synthetic
heatmap directly comparable to real sensor output.

%% TODO: Name the specific body model and physics engine used, and report
%% the parameter sampling ranges (height, weight, pose).

\subsection{Simulation Environment}
\label{sec:simulation-environment}

%%TODO: Describe the Pybullet simulation environment, including the software stack,
%% its softbody simulation capabilities, and any custom modifications made to support the seat contact simulation.

%%TODO: what are the consideration, limitation and features of the simulation environment, and 
%%how they affect the synthetic data generation process.

\subsection{Heatmap Rendering }
\label{sec:heatmap-rendering}
The rendering pipeline samples occupant shape, pose, and seating
configuration parameters, runs the physics-based contact simulation, and
rasterises the resulting per-cell contact pressure into the same heatmap
format used for real sensor data, allowing synthetic and real samples to
be interchanged during model training without format-specific handling.
%% TODO: Document the rendering pipeline's software stack e.


\subsection{Augmentation Strategy}
\label{sec:synthetic-augmentation}
To increase the diversity of the synthetic dataset beyond the sampled
body model parameters, additional augmentations are applied, including
random sensor noise injection, simulated calibration drift, and small
random seat-position offsets, so that models trained on synthetic data
are not overly reliant on idealised, noise-free heatmaps.

%% TODO: List exact augmentation and randomisation parameters and value ranges used.

\section{Sim2Real Validation}
\label{sec:sim2real-validation}

%%TODO: attach some figures showing the comparison between real and synthetic heatmaps, 
%% and the domain gap assessment results.


\section{Final data set pipeline }
\label{sec:final-dataset-composition}
Table~\ref{tab:dataset-composition} s.

\begin{table}[H]
\centering
\caption{Final dataset composition by split, showing real and synthetic
sample counts and class distribution.}`'
\label{tab:dataset-composition}
\begin{tabular}{@{}lrrrl@{}}
\toprule
\textbf{Split} & \textbf{Real Samples} & \textbf{Synthetic Samples} & \textbf{Total} & \textbf{Class Distribution} \\
\midrule
Training   & \todo{X} & \todo{X} & \todo{X} & \todo{e.g. 40\% adult / 25\% child / ...} \\
Validation & \todo{X} & \todo{X} & \todo{X} & \todo{e.g. 40\% adult / 25\% child / ...} \\
Test       & \todo{X} & \todo{X} & \todo{X} & \todo{e.g. 40\% adult / 25\% child / ...} \\
\midrule
\textbf{Total} & \todo{X} & \todo{X} & \todo{X} & -- \\
\bottomrule
\end{tabular}
\end{table}

%% TODO: Populate this table with the actual final sample counts and
%% class distribution percentages once data collection and generation
%% are complete. The full extended version of this table also appears
%% in Appendix~\ref{app:dataset-statistics}.
